\chapter{Introductions}
\hypertarget{00_introductions}{}\label{00_introductions}\index{Introductions@{Introductions}}
{\bfseries{Cracky\+OS scripting}} is a script interpreter designed to make it easier to program and design test systems while providing a range of capabilities and flexibility.

It supports procedural and functional programming. It accepts a script source file as input and executes the script. There is no need to recompile the script. The system runs the code in two passes\+: in the first pass, it translates the script source code into tokens; in the second pass, it runs the generated file. Precompiled files use the extension \"{}.\+dad\"{}, while source files use the extension \"{}.\+dat\"{}.

Because this is an interpreter-\/based OS, scripts run more slowly than compiled code performing the same task. The aim of the system is not speed, but flexibility and rapid development.

When debug logging is enabled, log messages are sent to the RS232 terminal. Each message uses a special code, which is documented in the error.\+h file. Debug logging can significantly slow down the system because of the limited speed of RS232. Use it only for the relevant parts of the code.

A GUI is also provided to make programming the system easier and help reduce coding errors. 